\documentclass[11pt,a4paper]{article}
\usepackage[T1]{fontenc}
\usepackage[utf8]{inputenc}
\usepackage{lmodern}
\usepackage[margin=2.5cm]{geometry}
\usepackage{microtype}
\usepackage{amsmath}
\usepackage{enumitem}
\usepackage[hidelinks]{hyperref}
\setlist{itemsep=1pt,topsep=3pt,parsep=0pt}

\begin{document}
\begin{center}
  {\Large\bfseries Sonifying the Heart: Inner Emotion Detection from ECG through Music}\\[6pt]
  Music Workshop Sprint 3: Individual Contribution\\[8pt]
  {\large Chervith Reddy} \quad \texttt{2024101076}
\end{center}
\noindent\rule{\textwidth}{0.4pt}

\section{My Role}
I built the first two stages of our pipeline: \textbf{ECG data} and \textbf{Heartbeats}. Every
clip that the music and emotion stages use comes from my part.

\section{What I Did}
\begin{itemize}
  \item \textbf{Data access.}
    \begin{itemize}
      \item Read \texttt{afdb} and \texttt{mitdb} straight from PhysioNet, streaming only the
        annotations and 30\,s of signal.
      \item Wrote a parallel downloader for the full databases, because PhysioNet limits each
        connection to about 20\,KB/s.
    \end{itemize}
  \item \textbf{Annotations.} Handled both formats: \texttt{afdb}'s separate beat (\texttt{.qrs})
    and rhythm (\texttt{.atr}) files, and \texttt{mitdb}'s single file that mixes beat labels and
    rhythm marks.
  \item \textbf{Heartbeat measures.} R--R intervals, heart rate and RMSSD for every clip.
  \item \textbf{Clip selection.} For each recording, one 30\,s normal clip (the steadiest of up to
    five windows) and one 30\,s AFib clip (the window with the most detected beats).
  \item \textbf{Records.}
    \begin{itemize}
      \item Chose five \texttt{afdb} recordings and \texttt{mitdb} 201 and 202.
      \item Excluded \texttt{mitdb} 219. Its ``normal'' stretch (RMSSD 270.6\,ms) is more
        irregular than its AFib (145.9\,ms), because of 133 blocked beats.
    \end{itemize}
  \item \textbf{Any ECG file.}
    \begin{itemize}
      \item Added XQRS heartbeat detection.
      \item On record 04043 it found 53 beats, a mean R--R of 559\,ms and an RMSSD of 14.2\,ms.
        The expert annotations give 54 beats, 559\,ms and 13.8\,ms.
    \end{itemize}
\end{itemize}

\section{Key Result}
\texttt{afdb}'s beat marks come from an automatic detector and jitter in places. Taking the first
30\,s made two normal clips look more irregular than AFib (record 04015: 199\,ms against 105\,ms).

After switching to the steadiest window, every recording separates clearly:
\begin{itemize}
  \item normal clips: RMSSD 7.6--64.1\,ms;
  \item AFib clips: RMSSD 99.9--194.3\,ms, 3--17 times the same recording's normal value.
\end{itemize}

\section{Next Steps}
\begin{itemize}
  \item Load the ECG of emotion-labelled datasets (WESAD, DREAMER, AMIGOS) to validate the
    emotion stage.
  \item Add the Sprint~2 items still open: SDNN, pNN50, Poincar\'e plots and rate-matched clips.
\end{itemize}

\end{document}
